\section{Preparation for proper pushforward}
\label{section-preparation-pushforward}

\begin{lemma}
\label{lemma-equal-dimension}
Let $(S, \delta)$ be as in Situation \ref{situation-setup}.
Let $X$, $Y$ be locally of finite type over $S$.
Let $f : X \to Y$ be a morphism.
Assume $X$, $Y$ integral and $\dim_\delta(X) = \dim_\delta(Y)$.
Then either $f(X)$ is contained in a proper closed subscheme
of $Y$, or $f$ is dominant and the extension of function fields
$R(X)/R(Y)$ is finite.
\end{lemma}

\begin{proof}
The closure $\overline{f(X)} \subset Y$ is irreducible as $X$
is irreducible (Topology, Lemmas
\ref{topology-lemma-image-irreducible-space} and
\ref{topology-lemma-irreducible}).
If $\overline{f(X)} \not = Y$, then we are done.
If $\overline{f(X)} = Y$, then $f$ is dominant and by
Morphisms,
Lemma \ref{morphisms-lemma-dominant-finite-number-irreducible-components}
we see that the generic point $\eta_Y$ of $Y$ is in the image of $f$.
Of course this implies that $f(\eta_X) = \eta_Y$, where $\eta_X \in X$
is the generic point of $X$. Since $\delta(\eta_X) = \delta(\eta_Y)$
we see that $R(Y) = \kappa(\eta_Y) \subset \kappa(\eta_X) = R(X)$
is an extension of transcendence degree $0$.
Hence $R(Y) \subset R(X)$ is a finite extension by
Morphisms, Lemma \ref{morphisms-lemma-finite-degree}
(which applies by
Morphisms, Lemma \ref{morphisms-lemma-permanence-finite-type}).
\end{proof}

\begin{lemma}
\label{lemma-quasi-compact-locally-finite}
Let $(S, \delta)$ be as in Situation \ref{situation-setup}.
Let $X$, $Y$ be locally of finite type over $S$.
Let $f : X \to Y$ be a morphism.
Assume $f$ is quasi-compact, and $\{Z_i\}_{i \in I}$ is a locally
finite collection of closed subsets of $X$.
Then $\{\overline{f(Z_i)}\}_{i \in I}$ is a locally finite
collection of closed subsets of $Y$.
\end{lemma}

\begin{proof}
Let $V \subset Y$ be a quasi-compact open subset.
Since $f$ is quasi-compact the open $f^{-1}(V)$ is
quasi-compact. Hence the set
$\{i \in I \mid Z_i \cap f^{-1}(V) \not = \emptyset \}$
is finite by a simple topological argument which we omit.
Since this is the same as the set
$$
\{i \in I \mid f(Z_i) \cap V \not = \emptyset \} =
\{i \in I \mid \overline{f(Z_i)} \cap V \not = \emptyset \}
$$
the lemma is proved.
\end{proof}









